\documentclass{article}
\usepackage[left=2cm, right=2cm, top=2cm, bottom=2cm]{geometry}
\usepackage{graphicx}
\usepackage{amsmath}
\usepackage{amssymb}
\usepackage{amsthm}
\usepackage{fancyhdr}
\usepackage{verbatim}
\usepackage{listings}
\usepackage{xcolor}
\usepackage{pdfpages}
\usepackage{cancel}
\usepackage{physics}

\lstset{
    language=C++,
    basicstyle=\ttfamily\footnotesize,
    keywordstyle=\color{blue},
    commentstyle=\color{green},
    stringstyle=\color{red},
    numbers=left,
    numberstyle=\tiny\color{gray},
    frame=single,
    breaklines=true,
    captionpos=b,
}


\title{MTH 553: Lecture \# 15}
\author{Cliff Sun}

\newtheorem{theorem}{Theorem}[section]
\newtheorem{lemma}[theorem]{Lemma}
\newtheorem{definition}[theorem]{Definition}
\newtheorem{conjecture}[theorem]{Conjecture}
\newtheorem{proposition}[theorem]{Proposition}
\newtheorem{corollary}[theorem]{Corollary}
\newtheorem{one minute paper}[theorem]{One Minute Paper}

\pagestyle{fancy}
\lhead{\textbf{\thepage}\ \ \nouppercase{\rightmark}}
\chead{MTH 553: Lecture \# 15}
\rhead{Cliff Sun}

\begin{document}

\maketitle

\section*{Conservation of energy}

The conservation of energy states that (in one dimension)

\begin{equation}
    0 = \frac{d}{dt}\frac{1}{2}\int_{-\infty}^{\infty}(u_t^2 + c^2 u_x^2)dx
\end{equation}

Here, 

\begin{equation}
    \frac{1}{2}\int_{\Omega}u_t^2 dx 
\end{equation}

Is the kinetic energy. Moreover, 

\begin{equation}
    \frac{1}{2}c^2 \int_{\Omega}u_x^2 dx 
\end{equation}

Is the potential enegry due to a stretching of the displacement force. Now, let's do higher dimensions.

\subsection*{Higher dimensions}

\begin{equation}
    u_{tt} = c^2 \nabla^2 u
\end{equation}

\textbf{Should learn this calculation.} Initial data is assumed to be compactly supported. Define energy 

\begin{equation}
    \epsilon(t) =  \frac{1}{2}\int_{\mathbb{R}^n}\left( u_t^2 + c^2|\nabla u |^2 \right)dx
\end{equation}

This is also simply just kinetic plus potential. Here, the claim is that the total energy is conserved. 

\begin{proof}
    First expand:
    \begin{equation}
        \epsilon'(t) = \int (u_t u_{tt} + c^2(\nabla u \cdot \nabla u_t) )dx
    \end{equation}
    Then, by the wave equation:
    \begin{equation}
        = c^2\int (u_t \nabla^2 u + \nabla u \cdot \nabla u_t)dx
    \end{equation}
    But this is merely just the product rule:
    \begin{equation}
        = c^2 \int \nabla \cdot \left(u_t \nabla u \right)dx
    \end{equation}
    Then, use the divergence theorem:
    \begin{equation}
        = 0 
    \end{equation}
    Since $u$ has compact support and therefore the boundary terms vanish at $\pm \infty$ go away. Therefore, energy is conserved. 
\end{proof}

\subsection*{Consequences}

\textbf{Uniqueness: }if $u_1$ and $u_2$ are solutions with data $g,h$, then $w = u_1 - u_2$ has zero initial data. Therefore,
\begin{equation}
    \epsilon_w(t) = \epsilon_w(0) = 0
\end{equation}

Therefore, $w_t = 0 \implies w = 0 \implies u_1 = u_2$.

\textbf{Continuous dependence: } Let $u_1$ have initial data $g_1,h_1$, $u_2$ similar and therefore $w$ with $g_1 - g_2$ and $h_1 - h_2$. Then 

\begin{equation}
    \epsilon_w(t) = \epsilon_w(0)
\end{equation}

So then

\begin{align*}
    &\int\left( w_t^2 + c^2|\nabla w|^2 \right)dx \\
    &= \int\left( (h_1 - h_2)^2 + c^2|\nabla (g_1 - g_2)|^2 \right)dx
\end{align*}

So then if $h_1 \approx h_2$ and $g_1 \approx g_2$. Then $w \approx 0$ and therefore $u_1 \approx u_2$. This is on the homework (probably just bounds). Here, the point here is that 
we can just use energy methods to yield results of which previously explicit formulas were needed. 

\subsection*{Domain of dependence by energy method}

Now, suppose a point $(x_0, t_0)$. With a domain of dependence in $\mathbb{R}^n$ centered in $x_0$. 

\begin{theorem}
    If $u$ is zero and $u_t$ is zero in $B(x_0, ct_0)$ (ball in $\mathbb{R}^n$ centered in $x_0$ with radius $ct$), then $u \equiv 0$ inside 
    the cone $C$. 
\end{theorem}

Given a ball $B(x_0, ct)$, then move the ball "up" in time to yield $B(x_0, c(t_0-t))$ where $u_t = 0$. Initial disturbances outside of $B(x_0,ct_0)$ cannot affect $u(x_0,t_0)$. So the wave travels at 
most speed $c$.

\begin{proof}
    Define a local energy (local with respect to slices of the cone $C$)
    \begin{equation}
        e(t) = \frac{1}{2}\int_{B(x_0, c(t_0 - t))}\left( u_t^2 + c^2 |\nabla u|^2 \right)dx \geq 0
    \end{equation}
    We prove that the local energy is \underline{dissipated} ($e'(t) \leq 0$). Indeed,
    \begin{align*}
        e'(t) &= \int_{B(x_0, c(t_0 - t))}\left[ u_t u_{tt} + c^2 \nabla u \cdot \nabla u_t \right]dx \\
        &=\int_{B(x_0, c(t_0 - t))}\left[ u_t \nabla \cdot \nabla u + c^2 \nabla u \cdot \nabla u_t \right]dx  + \frac{1}{2}\int_{\partial B(x_0, c(t_0 - t))}\left[ u_t^2 + c^2 |\nabla u|^2 \right]dS \cdot \frac{d}{dt}(c(t_0 -t))\\
    \end{align*}
    The last line is by the fundamental theorem of Calculus since 
    \begin{align*}
        &\frac{d}{dr}\int_{B^n(r)}f(x)dx \\
        &= \frac{d}{dr}\int_0^r \int_{S^{n-1}}f(\rho \xi)d\xi \rho^{n-1}d\rho \\
        &= \int_{S^{n-1}}f(r\xi)d\xi r^{n-1} \text{ by FTC} \\
        &= \int_{S^{n-1}(r)}fdS
    \end{align*}
    \begin{align*}
        e'(t) &= c^2 \int_{B(x_0, c(t_0 - t))}\left[ \nabla \cdot (u_t \nabla u) \right]dx - \frac{c}{2}\int_{\partial B}\left[ u_t^2 + c^2 |\nabla u|^2 \right]dS\\
        &= c\int_{\partial B}\left[ cu_t \nabla u \cdot \vec{v} - \frac{1}{2}u_t^2 - \frac{c^2}{2}|\nabla u|^2 \right]dS
    \end{align*}
    Note that this is less than or equal to $0$ since 
    \begin{equation}
        |cu_t\nabla u \cdot v| \leq |u_t| c|\nabla u| \leq \frac{1}{2}u_t^2 + \frac{1}{2}c^2 |\nabla u|^2 
    \end{equation}
    Since $2ab \leq a^2 + b^2$. But if $e(t) \leq e(0) = 0$, then $e(t) = 0$ which then implies that $u_t = 0$ and therefore $u = 0$. 
\end{proof}

\end{document}